% ============================================================
%  Chapter 26 (Coda): The Union
% ============================================================

\epigraph{%
  The two crowns are not two.\\
  They never were.%
}{}

\section{What Has Been Shown}
\label{sec:ch26-what}

We began with an ion.
A single charged molecule, $m/z = 299.0555$, moving through a mass spectrometer.
We described its journey in two languages: classical Newtonian mechanics and
quantum wave mechanics.
Both descriptions gave the same number at the detector.

We then asked: why do they agree?

The answer, worked out over twenty-four chapters and seven physical domains,
is this: because the ion's phase space is bounded.

Bounded phase space forces oscillation.
Oscillation has a discrete spectrum.
A discrete spectrum generates quantum numbers.
Quantum numbers label partition cells.
Partition cells have boundaries.
Boundaries carry charge.
Charged bounded oscillators have rest mass.
Rest mass requires Lorentz invariance.
Lorentz invariance with a finite propagation bound is special relativity.

None of these steps required an additional postulate.
They are the logical content of three words: \emph{systems are bounded}.

\section{The Two Crowns}
\label{sec:ch26-crowns}

The First Crown --- classical mechanics --- and the Second Crown ---
quantum mechanics --- have been shown to be two readings of the same text.
The text is the bounded phase space.

Classical mechanics reads the text as trajectories: positions, momenta, and
forces.
Quantum mechanics reads it as wave functions: eigenstates, operators, and
probabilities.
Both readings are correct.
Neither is more fundamental.
The Triple Equivalence $\Osc \cong \Cat \cong \Part$ is the formal proof
that the two readings are isomorphic.

\section{The Union}
\label{sec:ch26-theunion}

The union is not a compromise between the two crowns.
It is the recognition that they were never separate.

The mass spectrometer is a bounded physical system.
The ion inside it is a bounded physical system.
The GPU that simulates its trajectory is, by the Triple Equivalence,
a physical observation apparatus.
The spectral database is the partition hierarchy.
The purpose function is the analytical intent encoded as a partition filter.
The S-entropy triple $(\Sk, \St, \Se)$ is the epistemic state of the
instrument.

All of these are the same structure, seen from different vantage points.

The monograph ends where it began: with a single ion.
It has mass.
It has charge.
It oscillates in a bounded phase space.
It is measured to better than two parts per million on four different instruments.

It agrees with itself.

\section{The Five Unions, Revisited}
\label{sec:ch26-prologue-echo}

The prologue to this monograph described five unions.

Imhotep compounded knowledge under a unified administration and produced
stone that looked impossible to those outside his workshop.
The credibility inversion of equation~\eqref{eq:credibility-inversion}
has not disappeared: the framework in this monograph will read as
overclaiming to readers whose priors are shaped by the Standard Model.
That is the structure of the problem, not a refutation of the content.

The van Aken workshop composed knowledge across three generations and
produced a density of symbolic reference that no single mind could have
maintained.
The composition inflation of Chapter~\ref{ch:inflation} --- $d(d+1)^{n-1}$
distinguishable trajectories at partition depth $n$ --- is the same
principle: what appears impossible for one level of description is routine
at the next.

Newton unified terrestrial and celestial mechanics and left $G$ as a
measured constant.
Chapter~\ref{ch:gravity} shows it is a derived quantity --- the ratio of
partition volumes in a three-dimensional bounded space --- recoverable
from first principles to arbitrary precision.

Galileo worked at the boundary between ray optics and wave optics and
left their reconciliation to the future.
Chapters~\ref{ch:lorentz} and~\ref{ch:triple} show that wave and particle
are the oscillatory and partition readings of the same bounded phase space:
not reconciled by compromise, but unified by the Triple Equivalence
$\mathfrak{O} \cong \mathfrak{E} \cong \mathfrak{P}$.

The fifth union --- the one this monograph reports --- subsumes all four.

\section{What Remains}
\label{sec:ch26-remains}

The bounded phase space framework does not claim to be complete.
The numerical values of fundamental constants ($\alpha$, $m_p/m_e$,
$G m_e^2/(\hbar c)$) remain un-derived here.
The spatial dimension is an input.
The generations of quarks and leptons are not addressed.
The full apparatus of quantum field theory --- creation and annihilation operators,
gauge bosons, renormalisation --- remains to be connected rigorously.
General relativity in its full form (Einstein field equations) is not derived.

Each of these is an open question, not a refutation.
Each can in principle be addressed without introducing a new axiom.
The programme is not finished; it is begun.

\medskip
The axiom is stated.
The consequences are derived.
The predictions are checked.

\emph{What remains, after the Axiom is stated and its Consequences deduced, is
the physical world as it is.}

\vfill
\begin{flushright}
\rule{4cm}{0.4pt}\\
\textit{K.F.S.}
\end{flushright}
